\section{Introduction}\label{sec:intro}
% Sim version (track sim, 2026-11-13). Rewritten with the journal-paper-writing skill: OCAR, about 900 words,
% a preview that gives directions without numbers (CONTRACT Sec. 5.7). Every number of the old introduction
% is stated once in Sec. 3. Outline, sentence table and critiques: paper-rework/jpw-review/sim/intro-*.md and
% critiques/intro_critique_round*.md (revised after critic rounds 1 to 3).
% The sentence on the IQM Garnet QPU and its TODO(authors) on the coupler map are not in the sim version;
% hardware appears only as an outlook in Sec. 4 (CONTRACT Sec. 6).

Encoding quantum information in an error-correcting code is worth its overhead only below the pseudo-threshold, the physical error rate at which a logical gate and a physical gate fail equally often. For qubits, logical and physical gates have been compared on hardware. Logical randomized benchmarking (logical RB) in the $[[4,2,2]]$ error-detecting code found a lower infidelity for logical gates than for their physical counterparts \cite{Harper2019FTIBM}, and logical RB has since benchmarked the transversal Clifford gates of a color code \cite{Lacroix2025ColorCode}. Qutrits now call for the same comparison: a third level shortens circuits such as that of the generalized Toffoli gate \cite{Gokhale2019,Wang2020}, magic-state distillation improves with the local dimension \cite{CampbellPRX2012,CampbellPRL2014}, and superconducting and trapped-ion processors provide a third level natively \cite{Goss2024,Ringbauer2022}.
% TODO(authors): confirm what Refs. CampbellPRX2012 and CampbellPRL2014 show for the local dimension three. The former
% sentence of the paper said that Reed-Muller codes give magic-state distillation in all prime dimensions and that
% working with more than two levels improves its overhead. A critic who read the abstracts and passages of the first
% paper found the large gain in yield reported for five levels, and for qutrits a yield of the same order as the
% qubit protocol; the clause above therefore claims only the trend with the dimension.

% TODO(authors): confirm the gap statement below against the recent literature on qudit error correction
% (logical qudit memories and qutrit topological codes on hardware); add citations if a logical qutrit gate
% has been benchmarked elsewhere.
For qutrits only parts of this comparison exist. At the physical level, qutrit Clifford gates have been characterized by randomized benchmarking (RB) \cite{Morvan2021}. At the logical level, error thresholds of qudit topological codes have been estimated for a quantum memory under code-capacity noise \cite{DuclosCianci2013}. No logical qutrit gate has been benchmarked on hardware or under circuit-level noise, and the pseudo-threshold of such a gate is unknown.

A first answer can come from simulation, and an experiment can test it if the code is small enough for near-term processors and has shallow logical gates. One such code is the $[[5,1,2]]_3$ code, a surface-code patch of five data qutrits that needs at most four ancillas for syndrome extraction. Its logical Clifford gates are transversal in the sense of the folded surface code \cite{Moussa2016}: they consist of one-qutrit gates and of two-qutrit gates on one pair of data qutrits.

The price of this small size is the distance, $d = 2$: the code detects every single-qutrit error and corrects none. The error correction that standard logical RB applies after each gate \cite{Combes2017} must therefore be replaced by postselection, which accepts a shot only if its syndromes are trivial.

Besides the noise itself, we consider three factors that set where the pseudo-threshold of such a code lies. The first is the postselection schedule, that is, how often the syndromes enter the postselection. Undetected errors persist in syndrome subspaces and act on later gates, so logical dynamics can be non-Markovian even under Markovian physical noise \cite{Ziyad2025}. The less often the syndromes enter the postselection, the longer such errors persist and can combine with later ones, and the logical signal of the accepted shots then need not decay as the single exponential from which RB extracts an error rate. The second factor is the layout of the circuit that extracts the syndromes: each round of syndrome extraction, or check round, adds faults of its own, and the connectivity of the processor sets its length. The third is the local dimension, which changes both the code and the physical gates that its logical gates are compared with.

We therefore ask: below which two-qutrit gate error probability $p$ does a postselected logical Clifford gate of the $[[5,1,2]]_3$ code have a lower error rate than a physical single-qutrit Clifford gate, and how does this pseudo-threshold $p^\star$ depend on the postselection schedule, on the layout of the check round and on the local dimension? To answer this, we simulate logical randomized benchmarking with detection (LRB-D), in which a check round follows each gate of a random sequence of logical Clifford gates, and the sequence is then inverted gate by gate with a single check round at the end. We use two circuit-level noise models, qutrit forms of the superconducting-inspired model (SI1000) and the standard depolarizing model (SD6) \cite{GidneyHoneycomb2021}, both parameterized by $p$. We sample every circuit with the \texttt{sdim} stabilizer simulator \cite{sdim} through the exact detector error model of the circuit, and we compare the fitted logical error rate with that of physical single-qutrit RB under the same noise model at the same $p$, in which each Clifford gate is the same sequence of primitive gates as the logical one.

We vary all three factors. For the schedule, every check round enters the postselection, or only every second, third or fourth one; we also simulate circuits that contain no check rounds and are postselected on the final data readout alone (terminal postselection). For the layout, we compare two arrangements of the data and ancilla qutrits. In the first, every stabilizer has an ancilla next to all of its data qutrits, and the two-qutrit gates of the logical Cliffords need one long-range coupler. The second needs no such coupler and has a longer check round. The simulations give the gates on the long-range coupler the noise of any other two-qutrit gate, so this comparison isolates the check round. For the local dimension, we repeat the simulations for the qubit $[[5,1,2]]$ code.

The simulations show that a logical Clifford of the $[[5,1,2]]_3$ code has a lower error rate than a physical one below a pseudo-threshold. This holds among the accepted shots, whose fraction falls with the depth of the gate sequence and with $p$. When every check round enters the postselection, the logical error rate grows about as $p^2$ at small $p$ and the physical one as $p$. The pseudo-threshold falls as fewer check rounds enter the postselection, and under terminal postselection the accepted logical signal does not decay as a single exponential, so we extract no error rate from it. The layout with the longer check round has the lower pseudo-thresholds, by a larger fraction under SD6. The qubit code has lower pseudo-thresholds than the qutrit code in both layouts, under both noise models and for each schedule with check rounds. These lower values combine two differences at equal $p$: below the qutrit pseudo-thresholds, wherever the fits resolve it, the logical gates of the qubit code have the higher error rate, by a mechanism that remains open, and its physical gates have the lower one.

In the following sections we define the code, the LRB-D protocol, the layouts and the noise models and say how we estimate $p^\star$; we then present the results and state what qualifies them, what sets the pseudo-threshold and what follows for an experiment. The appendices hold the check-round circuits, the details and tables of the simulations and fits, and the representation theory of the twirled channel.
